\section{Elektrotechnische Grundlagen -- Formelsammlung}

Diese Section sammelt \textbf{elektrotechnische Kernformeln} für schnelle Rechnungen
in Leistungselektronik, Netzwerken, Maschinen und Feldern:
\textbf{Strom/Spannung/Leistung}, \textbf{RLC}, \textbf{Wechselstromrechnung},
\textbf{Magnetismus}, \textbf{Transformator}, \textbf{Drehstrom} und \textbf{Maschinen-Base}.
Sie ist als \textbf{Nachschlage-Section} aufgebaut.

% -------------------------------------------------
\subsection{Grundgrössen und Definitionen}

\subsubsection{Strom, Spannung, Leistung}

\[
\cbl{i(t) = \frac{dq(t)}{dt}}
\qquad
\cbl{u(t) = \frac{dw(t)}{dq(t)}}
\]

\[
\cbl{p(t) = u(t)\, i(t)}
\qquad
\cbl{w(t) = \int p(t)\,dt}
\]

\subsubsection{Passives Vorzeichenkonzept}

\[
\cbl{p>0 \Rightarrow \text{Verbraucherbetrieb}}
\qquad
\cbl{p<0 \Rightarrow \text{Generatorbetrieb}}
\]

% -------------------------------------------------
\subsection{Kirchhoff und Netzwerkanalyse}

\subsubsection{Knoten- und Maschensatz}

\[
\cbl{\sum i_k = 0 \quad \text{(KCL)}}
\qquad
\cbl{\sum u_k = 0 \quad \text{(KVL)}}
\]

\subsubsection{Leiterwerte und Widerstände}

\[
\cbl{R = \rho \frac{\ell}{A}}
\qquad
\cbl{G = \frac{1}{R}}
\]

\subsubsection{Serien- und Parallelschaltung}

\[
\cbl{R_{\mathrm{Serie}} = \sum R_k}
\qquad
\cbl{\frac{1}{R_{\mathrm{Par}}} = \sum \frac{1}{R_k}}
\]

% -------------------------------------------------
\subsection{Bauelementgleichungen (Zeitbereich)}

\subsubsection{Widerstand}

\[
\cbl{u_R(t) = R\, i(t)}
\qquad
\cbl{p_R(t) = R\, i^2(t) = \frac{u_R^2(t)}{R}}
\]

\subsubsection{Induktivität}

\[
\cbl{u_L(t) = L \frac{di(t)}{dt}}
\qquad
\cbl{i_L(t) \;\text{stetig}}
\]

Energie:
\[
\cbl{w_L = \frac{1}{2} L i^2}
\]

\subsubsection{Kapazität}

\[
\cbl{i_C(t) = C \frac{du(t)}{dt}}
\qquad
\cbl{u_C(t) \;\text{stetig}}
\]

Energie:
\[
\cbl{w_C = \frac{1}{2} C u^2}
\]

% -------------------------------------------------
\subsection{Wechselstromrechnung (Sinus, Phasor)}

\subsubsection{Sinusgrössen}

\[
\cbl{u(t) = \hat U \sin(\omega t + \varphi)}
\qquad
\cbl{i(t) = \hat I \sin(\omega t + \varphi_i)}
\]

\[
\cbl{\omega = 2\pi f}
\]

Effektivwerte:
\[
\cbl{U = \frac{\hat U}{\sqrt{2}}}
\qquad
\cbl{I = \frac{\hat I}{\sqrt{2}}}
\]

\subsubsection{Phasor-Notation}

\[
\cbl{\underline{U} = U \e^{\jimg \varphi}}
\qquad
\cbl{\underline{I} = I \e^{\jimg \varphi_i}}
\]

\subsubsection{Impedanzen}

\[
\cbl{Z_R = R}
\qquad
\cbl{Z_L = \jimg \omega L}
\qquad
\cbl{Z_C = \frac{1}{\jimg \omega C}}
\]

Admittanzen:
\[
\cbl{Y = \frac{1}{Z}}
\qquad
\cbl{Y_R = \frac{1}{R}}
\qquad
\cbl{Y_L = \frac{1}{\jimg \omega L}}
\qquad
\cbl{Y_C = \jimg \omega C}
\]

\subsubsection{Komplexe Leistung}

\[
\cbl{\underline{S} = \underline{U}\, \underline{I}^{\,\ast} = P + \jimg Q}
\]

\[
\cbl{P = U I \cos\varphi}
\qquad
\cbl{Q = U I \sin\varphi}
\qquad
\cbl{|S| = U I}
\]

Leistungsfaktor:
\[
\cbl{\lambda = \frac{P}{|S|} = \cos\varphi \quad \text{(Sinusfall)}}
\]

% -------------------------------------------------
\subsection{Nichtsinusförmige Grössen (RMS, Mittelwert)}

Mittelwert (über Periode $T$):
\[
\cbl{\overline{x} = \frac{1}{T}\int_0^T x(t)\,dt}
\]

Effektivwert:
\[
\cbl{X_{\mathrm{RMS}} = \sqrt{\frac{1}{T}\int_0^T x^2(t)\,dt}}
\]

Wirkleistung allgemein:
\[
\cbl{P = \frac{1}{T}\int_0^T u(t)i(t)\,dt}
\]

% -------------------------------------------------
\subsection{Magnetismus und Felder}

\subsubsection{Grundgrössen}

\[
\cbl{\vec{B} \;[\mathrm{T}], \quad \vec{H}\;[\mathrm{A/m}], \quad \Phi\;[\mathrm{Wb}]}
\]

\[
\cbl{\Phi = \int \vec{B}\cdot d\vec{A}}
\qquad
\cbl{B = \frac{\Phi}{A}}
\]

Materialgleichung (linear):
\[
\cbl{\vec{B} = \mu \vec{H}}
\qquad
\cbl{\mu = \mu_0 \mu_r}
\qquad
\cbl{\mu_0 = 4\pi \cdot 10^{-7}\,\mathrm{H/m}}
\]

\subsubsection{Ampèresches Gesetz}

\[
\cbl{\oint \vec{H}\cdot d\vec{\ell} = N i}
\]

Für magnetischen Kreis:
\[
\cbl{H \ell = N i \Rightarrow H = \frac{N i}{\ell}}
\]

\subsubsection{Magnetische Spannung, Durchflutung}

\[
\cbl{\Theta = N i \quad \text{(Durchflutung)}}
\qquad
\cbl{\mathcal{F}=\Theta}
\]

\subsubsection{Reluktanz und magnetischer Widerstand}

\[
\cbl{\mathcal{R} = \frac{\ell}{\mu A}}
\]

Magnetischer Ohmscher Zusammenhang:
\[
\cbl{\Phi = \frac{\Theta}{\mathcal{R}}}
\]

Reluktanzen in Serie/Parallel:
\[
\cbl{\mathcal{R}_{\mathrm{Serie}} = \sum \mathcal{R}_k}
\qquad
\cbl{\frac{1}{\mathcal{R}_{\mathrm{Par}}} = \sum \frac{1}{\mathcal{R}_k}}
\]

\subsubsection{Induktivität aus Magnetkreis}

\[
\cbl{L = \frac{N\Phi}{i}}
\qquad
\Rightarrow
\qquad
\cbl{L = \frac{N^2}{\mathcal{R}}}
\]

\subsubsection{Luftspalt (dominant)}

\[
\cbl{\mathcal{R}_\delta = \frac{\delta}{\mu_0 A}}
\]

Wenn Luftspalt dominiert:
\[
\cbl{L \approx \frac{N^2 \mu_0 A}{\delta}}
\]

\subsubsection{Energie und Kraft im Magnetfeld}

Energiedichte:
\[
\cbl{w_m = \frac{1}{2} \vec{B}\cdot \vec{H}}
\]

Gesamtenergie (linear, Induktivität):
\[
\cbl{W_m = \frac{1}{2} L i^2}
\]

Kraft im Luftspalt (vereinfachtes Modell):
\[
\cbl{F \approx \frac{B^2}{2\mu_0} A}
\]

\subsubsection{Induzierte Spannung (Faraday)}

\[
\cbl{u_{\mathrm{ind}}(t) = - N \frac{d\Phi(t)}{dt}}
\]

\[
\cbl{\nabla \times \vec{E} = -\frac{\partial \vec{B}}{\partial t}}
\]

% -------------------------------------------------
\subsection{Transformator (Ideal und Näherung)}

\subsubsection{Idealer Trafo}

Windungsverhältnis:
\[
\cbl{n = \frac{N_1}{N_2}}
\]

Spannungsübersetzung:
\[
\cbl{\frac{U_1}{U_2} = \frac{N_1}{N_2} = n}
\]

Stromübersetzung:
\[
\cbl{\frac{I_1}{I_2} = \frac{N_2}{N_1} = \frac{1}{n}}
\]

Leistung (ideal):
\[
\cbl{P_1 = P_2}
\]

Impedanzreflexion:
\[
\cbl{Z_1 = n^2 Z_2}
\]

\subsubsection{Induktivitäten und Kopplung}

\[
\cbl{L_1 = \frac{N_1^2}{\mathcal{R}}}
\qquad
\cbl{L_2 = \frac{N_2^2}{\mathcal{R}}}
\qquad
\cbl{M = k \sqrt{L_1 L_2}}
\]

% -------------------------------------------------
\subsection{Drehstrom (3~) -- Kernformeln}

\subsubsection{Spannungen (symmetrisch)}

\[
\cbl{u_a(t)=\hat U_{ph}\sin(\omega t)}
\]
\[
\cbl{u_b(t)=\hat U_{ph}\sin(\omega t-120^\circ)}
\]
\[
\cbl{u_c(t)=\hat U_{ph}\sin(\omega t-240^\circ)}
\]

Verhältnis Leiter--Phase:
\[
\cbl{U_{LL} = \sqrt{3}\,U_{ph}}
\qquad
\cbl{\hat U_{LL} = \sqrt{3}\,\hat U_{ph}}
\]

\subsubsection{Leistung im Drehstromsystem}

\[
\cbl{P = \sqrt{3}\,U_{LL}\,I_L\cos\varphi}
\]
\[
\cbl{Q = \sqrt{3}\,U_{LL}\,I_L\sin\varphi}
\]
\[
\cbl{S = \sqrt{3}\,U_{LL}\,I_L}
\]

\subsubsection{Stern/Dreieck}

Stern:
\[
\cbl{U_{ph}=\frac{U_{LL}}{\sqrt{3}}}
\qquad
\cbl{I_{ph}=I_L}
\]

Dreieck:
\[
\cbl{U_{ph}=U_{LL}}
\qquad
\cbl{I_{ph}=\frac{I_L}{\sqrt{3}}}
\]

% -------------------------------------------------
\subsection{Elektrische Maschinen -- Basisformeln}

\subsubsection{Elektrische und mechanische Winkelgeschwindigkeit}

\[
\cbl{\omega = 2\pi f}
\qquad
\cbl{n = \frac{\omega}{2\pi}\cdot 60 \;\; [\mathrm{rpm}]}
\]

\subsubsection{Synchronmaschine (Grundbeziehungen)}

Synchronwinkelgeschwindigkeit:
\[
\cbl{\omega_s = \frac{2\pi f}{p}}
\qquad
\cbl{n_s = \frac{60 f}{p}}
\]
($p$ = Polpaarzahl)

\subsubsection{Asynchronmaschine (Schlupf)}

\[
\cbl{s = \frac{n_s - n}{n_s}}
\qquad
\cbl{f_r = s f}
\]

\subsubsection{Gleichstrommaschine (Kernformeln)}

\[
\cbl{E_i = \omega \Psi}
\qquad
\cbl{M = i_a \Psi}
\qquad
\cbl{u_a = R_a i_a + E_i}
\]

% -------------------------------------------------
\subsection{Leitung, Feld und Lorentzkraft}

\subsubsection{Lorentzkraft}

\[
\cbl{\vec{F} = q(\vec{E} + \vec{v}\times \vec{B})}
\]

Leiterkraft (vereinfachtes Modell):
\[
\cbl{\vec{F} = I \, \vec{\ell} \times \vec{B}}
\]
Betrag:
\[
\cbl{F = B I \ell \sin\theta}
\]

\subsubsection{Leistungsfluss (Poynting)}

\[
\cbl{\vec{S} = \vec{E}\times \vec{H}}
\]

% -------------------------------------------------
\subsection{Typische Merksätze (ET-Kern)}

\begin{itemize}
    \item $u_L=L\,di/dt$ und $i_C=C\,du/dt$ sind die wichtigsten Zeitableitungsbeziehungen.
    \item $W_L=\tfrac{1}{2}Li^2$ und $W_C=\tfrac{1}{2}Cu^2$ sind prüfungsrelevant.
    \item Magnetkreis: $\Phi=\Theta/\mathcal{R}$ und $L=N^2/\mathcal{R}$.
    \item Drehstromleistung: $P=\sqrt{3}U_{LL}I_L\cos\varphi$.
\end{itemize}

\subsection{Typische Fehlerquellen}

\begin{itemize}
    \item Leiter- und Phasenspannung verwechselt ($\sqrt{3}$-Faktor).
    \item Vorzeichen bei Induktion ($u=-N d\Phi/dt$) ignoriert.
    \item Luftspalt-Reluktanz unterschätzt.
    \item RMS mit Scheitelwert verwechselt.
\end{itemize}
